\scanpage{326}
\mspage{165}

strict geometric Teichmüller (once the ``origin'' element $\Sigma$ has been chosen). More \unsure{intrinsically}, we \unsure{will have}:
\[
(17)\qquad
\begin{gathered}
\text{Set of strict geometric}\\
\text{Teichmüller subgroups}\\
\text{of }\widehat{\widehat{\mathcal{T}}}(\pi)
\end{gathered}
\unsure{\xleftarrow{\sim}|\simeq}
\operatorname{Isomext}_{\unsure{\mathrm{lac}}}
(\widehat{\Pi}_{g,\nu},\pi)/\unsure{\mathcal{N}_{g,\nu}}
\]
where we set, as is only right
\[
(18)\qquad
\mathcal{N}_{g,\nu}
=\operatorname{Norm}_{\widehat{\widehat{\mathcal{T}}}_{g,\nu}}
\bigl(\widehat{\mathcal{T}}_{g,\nu}^{+}\bigr)
\]
(perhaps equal to all of $\widehat{\widehat{\mathcal{T}}}_{g,\nu}$!)

Let us now consider the \uline{modular topos} \uline{over} $\mathbb{Q}$ of algebraic curves of type $(g,\nu)$, denoted $M_{g,\nu\,\mathbb{Q}}$, or simply $M_{g,\nu}$. If we choose a universal covering $\widetilde{M}_{g,\nu}$ of $M_{g,\nu\,\mathbb{Q}}$ (hence a universal covering of $\unsure{\operatorname{Spec}}\mathbb{Q}$ i.e. an algebraic closure $\overline{\mathbb{Q}}$ of $\mathbb{Q}$) we can \unsure{specify} the $\pi_1(M_{g,\nu\,\mathbb{Q}})$ as the group of $M_{g,\nu\,\mathbb{Q}}$-\unsure{automorphisms} of this \unsure{covering|covering}; \deleted{\unsure{and of it} \unsure{make} \illegible{1}} \inserted{call it the} \uline{arithmetic Teichmüller group} of type $(g,\nu)$ (relative to the choice of $\widetilde{M}_{g,\nu\,\mathbb{Q}}$). We will thus have an exact sequence

\mspage{325}

\scanpage{327}

\[
(19)\qquad 1\longrightarrow \pi_1(M_{g,\nu\,\overline{Q}})
\longrightarrow \pi_1(M_{g,\nu})
\longrightarrow \operatorname{Gal}(\overline{Q}/Q)\longrightarrow 1
\]
\deleted{\(1\longrightarrow \unsure{\widehat{\mathcal{T}}_{g,\nu}}\longrightarrow \unsure{\Gamma}\)}

(\unsure{where} $\overline{Q}$ and the base points are \unsure{made explicit} as has been said) Let us note that by the \unsure{transcendental} theory of Teichmüller, we have an isomorphism (defined modulo \unsure{operation} \unsure{induced} by $\pi_1(M_{g,\nu})$ on $\pi_1(M_{g,\nu,\overline{Q}})$)
\[
(20)\qquad \pi_1(\unsure{M_{g,\nu,\overline{Q}}})\simeq\widehat{\mathcal{T}}_{g,\nu}^{+}
\qquad (\text{profinite compactification of }\mathcal{T}_{g,\nu}^{+})
\]

Let us on the other hand consider the \unsure{scheme.} $\unsure{U|\mathcal{U}}_{g,\nu}$\inserted{ over $M_{g,\nu}$}, ``universal'' curve of type $(g,\nu)$

We thus have, on choosing a universal covering $\widetilde{\unsure{U|\mathcal{U}}}_{g,\nu}$ of it \uline{\unsure{over}} the chosen universal covering $\unsure{\widetilde{M}_{g,\nu}^{\natural}}$ of $M_{g,\nu}$, an exact sequence
\diagram{
\begin{tikzcd}[column sep=large,row sep=large,ampersand replacement=\&]
(21) \& 1 \arrow[r] \& \pi \arrow[r] \arrow[dd,"\Vert" description] \& \pi_1(\unsure{\overset{\deleted{\sim}}{\unsure{U|\mathcal{U}}}_{g,\nu}}) \arrow[r] \arrow[dd] \& \pi_1(M_{g,\nu}) \arrow[r] \arrow[dd] \& 1 \\
\text{which maps to:} \& \& \& \& \& \\
(22) \& 1 \arrow[r] \& \pi \arrow[r] \& \widehat{\widehat{\unsure{\mathfrak{S}|\mathcal{G}}}}(\pi) \arrow[r] \& \widehat{\widehat{\mathcal{T}}}(\pi) \arrow[r] \& 1
\end{tikzcd}
}
and the same for $\unsure{U|\mathcal{U}}_{g,\nu\,\overline{Q}}\longrightarrow M_{g,\nu\,\unsure{\overline{Q}|Q}}$

\mspage{326}

\scanpage{328}
\mspage{\unsure{166|168}}

\[
\unsure{(\overline{21})|(\overline{21}')}\qquad
1\longrightarrow\pi\longrightarrow\pi_1(\unsure{U|\mathcal{U}}_{g,\nu\,\overline{Q}})
\longrightarrow\pi_1(M_{g,\nu\,\overline{Q}})\longrightarrow1
\]
\deleted{\(\unsure{(\overline{22})}\qquad 1\longrightarrow\pi\longrightarrow\)}

We of course have an \unsure{isomorphism} of groups with \unsure{loops|ends}
\[
(23)\qquad \pi\simeq\widehat{\Pi}_{g,\nu}
\]
whose indeterminacy we would like to determine precisely

Let $M_{g,\nu\,C}$ denote the complex modular Teichmüller \deleted{\illegible{1}}\inserted{topos}, defined using the \deleted{differentiable}\inserted{$C^\infty$} surface $\unsure{U|\mathcal{U}}_{g,\nu}$, it is equipped with the canonical universal covering $\widetilde{M}_{g,\nu\,C}$ (\S~21)
\margin{$\widetilde{M}_{g,\nu\,C}=E_g/\unsure{A_{g,\nu}^{0}|A_{g,\nu}^{\circ}}$}
with group \unsure{of automorphisms} $\mathcal{T}_{g,\nu}^{+}$ precisely \unsure{--|...}
$\unsure{U|\mathcal{U}}_{g,\nu\,C}$ being itself also equipped with a universal covering \unsure{over} the preceding one,
$\widetilde{\unsure{U|\mathcal{U}}}_{g,\nu\,C}=E_g/\unsure{A_{g,\nu+1}^{0}|A_{g,\nu+1}^{\circ}}$.
This gives rise to the exact sequence of discrete groups (in the general topology context)
\diagram{
\begin{tikzcd}[column sep=large,row sep=large,ampersand replacement=\&]
(24) \& 1 \arrow[r] \& \Pi_{g,\nu} \arrow[r] \arrow[dd] \& \unsure{\mathfrak{S}|\mathcal{G}}_{g,\nu}^{+} \arrow[r] \arrow[d,equal] \& \mathcal{T}_{g,\nu}^{+} \arrow[r] \arrow[d,equal] \& 1 \\
\& \& \& \pi_1(\unsure{U|\mathcal{U}}_{g,\nu\,C}) \arrow[d] \& \pi_1(M_{g,\nu\,C}) \arrow[d] \& \\
\text{(mapping to} \& 1 \arrow[r] \& \Pi_{g,\nu} \arrow[r] \& \operatorname{Aut}_{\mathrm{lac}}\deleted{\unsure{\mathrm{ext}}}(\Pi_{g,\nu}) \arrow[r] \arrow[d,equal] \& \operatorname{Aut\,ext}_{\mathrm{lac}}(\Pi_{g,\nu}) \arrow[r] \arrow[d,equal] \& 1 \\
\& \& \& \unsure{\mathfrak{S}|\mathcal{G}}_{g,\nu} \& \mathcal{T}_{g,\nu} \&
\end{tikzcd}
}
by an \unsure{isomorphism|homomorphism} of extension \unsure{structures}), which by passage to the profinite completions gives the analogous exact sequence for the

\mspage{327}

\scanpage{329}

algebraic multiplicities over $\mathbb C$

\diagram{\begin{tikzcd}[ampersand replacement=\&,column sep=1.3em,row sep=1.2em]
(25) \& 1 \arrow[r] \&
\widehat{\Pi}_{g,\nu} \arrow[r] \arrow[d,equal] \&
\deleted{\unsure{\widehat{\Pi}}}\,\widehat{\mathcal G}_{g,\nu}
  \arrow[r] \arrow[d,very thick] \&
\widehat{\mathcal T}_{g,\nu}^{+} \arrow[r]
  \arrow[d,phantom,"{\rotatebox{-90}{$\unsure{\hookrightarrow}$}}" description] \& 1 \\
\& \& \inserted{\widehat{\Pi}_{g,\nu}} \arrow[r] \&
\inserted{\operatorname{Aut}\deleted{\unsure{\mathrm{ext}}}_{\mathrm{lac}}(\widehat{\Pi}_{g,\nu})}
  \arrow[r] \&
\inserted{\operatorname{Autext}_{\mathrm{lac}}(\widehat{\Pi}_{g,\nu})}
  \& \\[-0.7em]
\& \& \& |[xshift=-2.2em]| {\inserted{\raisebox{1ex}{\rotatebox[origin=c]{-60}{$=$}}\,\unsure{\widehat{\widehat{\mathcal G}}}_{g,\nu}}} \&
|[xshift=-2.2em]| {\inserted{\raisebox{1ex}{\rotatebox[origin=c]{-60}{$=$}}\,\widehat{\widehat{\mathcal T}}_{g,\nu}}} \&
\end{tikzcd}}

This situation moreover \unsure{arises} \unsure{canonically} from an analogous situation over $\overline{\mathbb Q}_0=\text{algebraic closure of }\mathbb Q\text{ in }\mathbb C$, the choice of $\widetilde{\unsure{U|\mathcal{U}}}_{g,\nu\mathbb C}$ defining a $\widetilde{\unsure{U|\mathcal{U}}}_{g,\nu\overline{\mathbb Q}_0}$, hence a $\widetilde{M}_{g,\nu\overline{\mathbb Q}_0}$.

To define an isomorphism between this sequence (25) and the exact sequence $(\overline{21})$ \deleted{one must} (\unsure{equipped} with \unsure{$2$ homomorphisms|its homomorphism} to (22)) one must

$1^{\circ})$ choose an isomorphism $\overline{\mathbb Q}\simeq\overline{\mathbb Q}_0$ --- the indeterminacy is in $\mathrm{Gal}(\overline{\mathbb Q}/\mathbb Q)$ ; which makes it possible to identify $\overline{\mathbb Q}$ and $\overline{\mathbb Q}_0$ ;

$2^{\circ})$ Choose an isomorphism (\deleted{\unsure{between}} $\overline{\mathbb Q}=\overline{\mathbb Q}_0$) of $(\widetilde{\unsure{U|\mathcal{U}}}_{g,\nu})_0$ (constructed by \unsure{transcendental|transcendental} means over $\mathbb C$ and descended to $\overline{\mathbb Q}_0$) with $\widetilde{\unsure{U|\mathcal{U}}}_{g,\nu\overline{\mathbb Q}}$ --- the indeterminacy is in $\pi_1(M_{g,\nu\overline{\mathbb Q}})$.

In summary, we have a class of isomorphisms of \unsure{(24)|(21)}\inserted{\unsure{2|?}} and (25), defined modulo automorphisms \inserted{\deleted{\illegible{1}} \illegible{1} \unsure{once we have chosen} $\overline{\mathbb Q}\simeq\overline{\mathbb Q}_0$} inner in \inserted{the \unsure{\uline{profinite}} group} $\pi_1(\unsure{U|\mathcal{U}}_{g,\nu})$. (\unsure{These|The} \inserted{\unsure{\uline{corresponding}}} isomorphisms transform the \deleted{structure} \unsure{Standard|standard} \uline{oriented} \uline{class of \unsure{discretification}} of $\widehat{\Pi}_{g,\nu}$ into \deleted{\unsure{a class}} \inserted{exactly} \uline{one} class of oriented discretifications of $\pi$.

\mspage{328}

\scanpage{330}
\mspage{167}

\deleted{Thus $\pi$ of (21) is \unsure{canonically} equipped with an oriented \deleted{\unsure{class of}} \inserted{\unsure{quasi}} discretification \inserted{\unsure{$\pi_0^{\natural}$}}, and in the homomorphism from $(\overline{21})$ to (22), we find that $\pi_1(\unsure{U|\mathcal{U}}_{g,\nu\overline{\mathbb Q}})$ is identified with $\widehat{\unsure{\mathfrak{S}|\mathcal{G}}}(\unsure{\pi_0^{\natural}})$, and $\pi_1(M_{g,\nu})$ with $\widehat{\mathcal T}(\unsure{\pi_0^{\natural}})$ \unsure{$\longrightarrow$}}

Changing the choice of the isomorphism to its complex conjugate, we find the opposite oriented \inserted{\unsure{quasi-}} \unsure{discretification}.

We see on the other hand that the image of $\pi_1(M_{g,\nu\overline{\mathbb Q}})\unsure{\subset}\pi_1(M_{g,\nu})$ in $\widehat{\widehat{\mathcal T}}(\pi)$ is none other than the $\widehat{\mathcal T}^{+}$ \unsure{corresponding|corresponding} to an \inserted{arbitrary} one of the chosen classes of discrete quasi-rigidifications --- the subgroup does not depend on the choice of this class --- which merely expresses that the image of $\pi_1(M_{g,\nu\overline{\mathbb Q}})$ in $\widehat{\widehat{\mathcal T}}(\pi)$, \unsure{considered|considered} as a subgroup of the image of the larger group $\pi_1(M_{g,\nu})$, is \unsure{normal|normal} in it --- which follows from the fact that $\pi_1(M_{g,\nu\overline{\mathbb Q}})$ is \unsure{normal} in $\pi_1(M_{g,\nu\unsure{\mathbb Q|\overline{\mathbb Q}}})$. Thus $\pi$

\mspage{329}